\subsection{Cross-Resolution Compactness Stability}

Table~\ref{tab:pp-by-state-resolution} presents mean Polsby-Popper scores
for the five focus states at three spatial resolutions.

\begin{table}[ht]
\centering
\caption{Mean Polsby-Popper compactness by state and resolution.
  All results single-run, seed 42.\textsuperscript{\dag}}
\label{tab:pp-by-state-resolution}
\begin{tabular}{lrrrr}
\toprule
State & County PP & Tract PP & Block-Group PP & RSI \\
\midrule
NC  ($k=14$) & 0.168 & 0.171 & 0.173 & 0.030 \\
TX  ($k=38$) & 0.194 & 0.198 & 0.200 & 0.030 \\
WI  ($k=8$)  & 0.201 & 0.203 & 0.205 & 0.020 \\
FL  ($k=28$) & 0.157 & 0.160 & 0.161 & 0.025 \\
VA  ($k=11$) & 0.189 & 0.191 & 0.193 & 0.021 \\
\midrule
\multicolumn{5}{l}{\footnotesize RSI = (max $-$ min) / mean across resolutions.
  Lower is more stable.} \\
\bottomrule
\end{tabular}
\end{table}

For all five states, PP increases slightly from county to block-group
resolution.  This is expected: finer units allow district boundaries to
follow physical geographic features more precisely, reducing the jagged
perimeter introduced when large county blocks straddle natural boundaries.
The \emph{magnitude} of the change is modest---between 2.0\% (WI) and 3.0\%
(NC, TX) in RSI terms.

\paragraph{Monotonic pattern.}
PP increases monotonically with resolution in all five states.  The
county-to-tract improvement ($\Delta PP \approx 0.002$--$0.004$) is roughly
equal in magnitude to the tract-to-block-group improvement
($\Delta PP \approx 0.001$--$0.002$).  This log-linear pattern is consistent
with the theoretical prediction that each tenfold increase in unit density
allows approximately equal perimeter refinement.

\subsection{50-State Summary}

Table~\ref{tab:50state-summary} reports the RSI distribution across all
50 states for the tract-to-block-group comparison (the most policy-relevant
comparison, since tracts are the research standard and block groups are the
finest resolution readily available for full 50-state analysis).

\begin{table}[ht]
\centering
\caption{Distribution of Resolution Stability Index (RSI) across 50 states,
  tract vs.\ block-group resolution.\textsuperscript{\dag}}
\label{tab:50state-summary}
\begin{tabular}{lr}
\toprule
RSI range & Number of states \\
\midrule
$< 1\%$                & 19 \\
$1\%$--$2\%$           & 29 \\
$2\%$--$3\%$           & 2  \\
$> 3\%$                & 0  \\
\midrule
\textbf{Total $<2\%$}  & \textbf{48} \\
\bottomrule
\end{tabular}
\end{table}

Forty-eight of 50 states exhibit RSI $< 2\%$, meeting the primary stability
criterion.\textsuperscript{\dag}  The two exceptions are Idaho (RSI = 2.3\%)
and Wyoming (RSI = 2.6\%), both of which have $k=2$ congressional districts.
At $k=2$ the bisect algorithm makes a single graph cut; small differences in
node granularity can shift which cut is optimal, producing non-trivial changes
in the resulting district shape.  This is a small-$k$ edge case rather than a
MAUP effect, as discussed in Section~5.

\subsection{D-Seat Stability Across Resolutions}

Table~\ref{tab:dseat-stability} reports D-seat counts at each resolution
level for the five focus states.

\begin{table}[ht]
\centering
\caption{Democratic seat count by state and resolution (2020 presidential
  benchmark).\textsuperscript{\dag}}
\label{tab:dseat-stability}
\begin{tabular}{lrrrr}
\toprule
State & $k$ & County D-seats & Tract D-seats & BG D-seats \\
\midrule
NC & 14 & 7 & 7 & 7 \\
TX & 38 & 13 & 13 & 13 \\
WI &  8 &  3 &  3 &  3 \\
FL & 28 & 10 & 10 & 10 \\
VA & 11 &  6 &  6 &  6 \\
\bottomrule
\end{tabular}
\end{table}

D-seat counts are identical across all three resolutions for all five focus
states.\textsuperscript{\dag}  This result is stronger than the compactness
finding: while PP shows a modest upward trend with resolution, the
seat-prediction metric is perfectly stable.  The result reflects that
partisan sorting operates at a geographic scale coarser than the unit size
tested: Democratic and Republican precincts cluster at the county and sub-county
level, not at the block-group level, so changes in unit granularity do not
move seats.

\subsection{Compactness Variance Decomposition}

To assess the relative contribution of resolution vs.\ geographic
factors to PP variation, we decompose total PP variance into
between-state and within-state (across-resolution) components:

\begin{align}
  \sigma^2_{\text{total}} &= \sigma^2_{\text{between-state}} +
                             \sigma^2_{\text{within-state (resolution)}} \\
  &\approx 0.0028 + 0.000006
\end{align}

Between-state variance dominates, accounting for 99.8\% of total variation.
Resolution contributes only 0.2\% of total PP variance.\textsuperscript{\dag}
This decomposition confirms that geographic characteristics (state shape,
coastlines, mountains) are the primary determinants of compactness, while
resolution choice is a negligible second-order factor.

\subsection{Outlier Analysis: Idaho and Wyoming}

The two states exceeding the 2\% RSI threshold are Idaho (RSI = 2.3\%) and
Wyoming (RSI = 2.6\%).  Both have $k=2$ congressional districts.  At $k=2$,
the bisect algorithm finds a single partition of the state into two equal-
population halves.  The county-level graph for Idaho has 44 nodes; the
optimal cut is very constrained and tends to place the boundary along the
Snake River corridor.  The block-group graph has $\sim$770 nodes, giving the
algorithm more flexibility to find an alternative cut.  The alternative cut
has similar population balance but a slightly different perimeter, shifting PP
by $\sim$0.005.  This is a small-$k$ sensitivity, not MAUP: the same
phenomenon would appear if we doubled the unit count further.
